\chapter{The action of the diagonal on boundary operators}
\label{ch:da-action}

A diagonal factorization contains two copies of the coordinates.
To act on a boundary, its second copy is identified with the
boundary coordinates. The two potential terms then cancel.
This produces a new factorization, but cancellation of the potential
does not yet identify its states with the original boundary.
That identification requires maps in both directions and a homotopy
between their composite and the identity.

Polynomial division supplies all three maps. Their failure to
preserve operator products strictly is itself controlled by the
same homotopy. The resulting comparison retains every ordered
product through explicit higher components.

\section{Eliminating one coordinate}
\label{sec:da-one-variable}

Let \(k\) be a field of characteristic zero, let \(W\in k[z]\), and
fix a scalar \(c\in k\). Let \(V\) be a finite-dimensional
\(k\)-vector space with a basis divided into even and odd vectors.
Put \(E(z)=k[z]\otimes_kV\), and let \(D(z)\) be an odd polynomial
matrix on this basis satisfying
\[
 D(z)^2=(W(z)-c)\,1_E.
\]
Thus \(E(z)\) is a finite free matrix factorization, with map
differential as in Section~\ref{sec:pb-boundary-differential}.
Introduce independent variables \(x,y\) and put
\[
 R=k[x],\quad S=k[x,y],\quad r=x-y,\quad
 g=\frac{W(x)-W(y)}{r},\quad D_x=D(x),\quad D_y=D(y).
\]
All quotients by \(r\) below are polynomial divided differences.
Their numerators vanish when \(y=x\); division by the monic
polynomial \(x-y\) therefore has zero remainder.

The diagonal factorization has an even basis vector \(1\), an odd
basis vector \(\theta\), and differential sending \(1\) to
\(g\theta\) and \(\theta\) to \(r\).
Its tensor product with \(E(y)\) is
\[
 M=S\otimes_k(V\oplus\theta V).
\]
On a homogeneous tensor the differential is
\(d_\Delta\otimes1+(-1)^{\text{parity in }\Delta}\otimes D_y\).
Writing a vector as \(v+\theta w\), this differential is
\begin{equation}
 Q(v+\theta w)
   =D_yv+rw+\theta(gv-D_yw).
 \label{eq:da-Q}
\end{equation}
The signs cancel the cross terms, giving
\[
 Q^2=(W(x)-W(y)+W(y)-c)\,1_M=(W(x)-c)\,1_M.
\]
When regarded only as an \(R\)-module, \(M\) is free on the
vectors \(y^j e\) and \(y^j\theta e\), for basis vectors \(e\)
and integers \(j\geq0\). Its rank over \(R\) is infinite when
\(E\ne0\).
Every polynomial vector is nevertheless a finite sum of these
vectors. This is the module obtained by forgetting the second
coordinate in the tensor product.

The intended finite replacement is \(N=E(x)\), with differential
\(D_x\). For a polynomial vector \(v(x,y)\), write
\[
 v_0=v(x,x),\qquad
 T(v)=\frac{v-v_0}{r},\qquad
 A=\frac{D_x-D_y}{r}.
\]
Define even maps
\(\pi:M\to N\), \(i:N\to M\), and an odd map \(h:M\to M\) by
\begin{equation}
 \pi(v+\theta w)=v_0,\qquad
 i(e)=e+\theta Ae,\qquad
 h(v+\theta w)=\theta T(v).
 \label{eq:da-contraction-maps}
\end{equation}
All are \(R\)-linear. They are not required to be \(S\)-linear:
the variable \(y\) has been eliminated in \(N\).

\begin{theorem}[The one-coordinate contraction]
\label{thm:da-one-contraction}
The maps in \eqref{eq:da-contraction-maps} satisfy
\begin{gather}
 \pi Q=D_x\pi,\qquad Qi=iD_x,\qquad \pi i=1_N,
 \label{eq:da-chain-maps}\\
 Qh+hQ=1_M-i\pi,\qquad
 h^2=0,\qquad \pi h=0,\qquad hi=0.
 \label{eq:da-homotopy}
\end{gather}
Hence \(\pi\) identifies the action of the diagonal with the
original boundary in the homotopy category.
No assumption on the critical points of \(W\) is needed.
\end{theorem}
\begin{proof}
The first component of \(Q(v+\theta w)\), evaluated at \(y=x\),
is \(D_xv_0\), which proves the equation for \(\pi\).
The equation for \(i\) reduces to two polynomial identities:
\[
 D_y+rA=D_x,\qquad D_yA+AD_x=g\,1_E.
\]
The first is the definition of \(A\). For the second,
\[
 D_yA+AD_x
 =\frac{D_yD_x-D_y^2+D_x^2-D_yD_x}{r}
 =\frac{W(x)-W(y)}r\,1_E .
\]
No commutation of \(D_x\) and \(D_y\) was used.

For the homotopy identity, \(Qh\) has components
\((v-v_0,-D_yT(v))\), while \(hQ\) has components
\((0,T(D_yv)+w)\).
The divided-difference identity
\[
 T(D_yv)=D_yT(v)-Av_0
\]
follows by writing \(v=v_0+rT(v)\) and using \(D_y-D_x=-rA\).
Adding the two components gives
\((v-v_0,w-Av_0)=(1_M-i\pi)(v+\theta w)\).
Finally, \(h\) kills the entire \(\theta E\) summand and vanishes
on vectors whose first component depends only on \(x\).
These observations prove the three side conditions, and
\(\pi i=1_N\) follows directly.
\end{proof}

The maps compare factorizations with the same scalar square.
They do not regard \(D_x\) or \(Q\) as square-zero differentials.
For a homogeneous map \(B\) between such factorizations, the
differential is the signed commutator
\(\partial B=d_{\mathrm{target}}B-(-1)^{|B|}Bd_{\mathrm{source}}\).
It squares to zero because the two scalar squares agree.
Thus \eqref{eq:da-homotopy} is an ordinary homotopy equation in
their map complex.

\section{Naturality and several coordinates}
\label{sec:da-several}

For a homogeneous polynomial map \(f:E\to F\) of parity \(p\),
tensoring with the diagonal gives the map
\[
 \mathcal T(f)(v+\theta w)
   =f(y)v+(-1)^p\theta f(y)w.
\]
The tensor signs show, by expansion of
\eqref{eq:da-Q}, that
\(\partial\mathcal T(f)=\mathcal T(\partial f)\).
They also give
\(\mathcal T(f_2f_1)=\mathcal T(f_2)\mathcal T(f_1)\).
The operation \(\mathcal T\) is therefore a differential graded
functor into factorizations on the polynomial pushforward modules.
Evaluation is strictly natural:
\begin{equation}
 \pi_F\mathcal T(f)=f(x)\pi_E
 \label{eq:da-natural}
\end{equation}
for every homogeneous map, whether closed or not.
For closed even maps it follows that the isomorphisms represented
by \(\pi_E\) are natural in the homotopy category.
The explicit inclusions \(i_E\) need not be strictly natural.

Let now \(n\geq1\), \(x=(x_1,\ldots,x_n)\), and
\(y=(y_1,\ldots,y_n)\).
Put \(R=k[x]\), \(S=k[x,y]\), and suppose
\(D(y)^2=(W(y)-c)\,1_E\).
For \(1\leq j\leq n\), set
\[
 r_j=x_j-y_j,\qquad
 g_j=
 \frac{
 W(x_1,\ldots,x_j,y_{j+1},\ldots,y_n)
 -W(x_1,\ldots,x_{j-1},y_j,\ldots,y_n)}
 {r_j}.
\]
The sum \(\sum_jr_jg_j=W(x)-W(y)\) telescopes.

The exterior module on \(\theta_1,\ldots,\theta_n\) has basis
the ordered products of distinct symbols, with
\(\theta_i^2=0\) and \(\theta_i\theta_j=-\theta_j\theta_i\).
Each \(\theta_i\) has odd parity. Let \(\epsilon_i\) denote left
multiplication by \(\theta_i\), and let \(\iota_i\) delete a
factor \(\theta_i\) with sign \((-1)^{j-1}\) when it occupies
position \(j\) in a product. It is zero on a product not containing
\(\theta_i\). These definitions imply
\[
 \iota_i\epsilon_j+\epsilon_j\iota_i=\delta_{ij},\qquad
 \epsilon_i\epsilon_j+\epsilon_j\epsilon_i=0,\qquad
 \iota_i\iota_j+\iota_j\iota_i=0.
\]
Here \(\delta_{ij}\) is \(1\) for \(i=j\) and \(0\) otherwise.
Consequently
\[
 D_\Delta=\sum_{j=1}^n(r_j\iota_j+g_j\epsilon_j),
 \qquad D_\Delta^2=(W(x)-W(y))\,1.
\]
On \(M=\Delta_W\otimes_{k[y]}E(y)\), put
\[
 Q(\omega\otimes e)
   =D_\Delta\omega\otimes e+(-1)^{|\omega|}\omega\otimes D(y)e.
\]
This has square \(W(x)-c\).
The module has finite rank over \(S\); for \(E\ne0\) its rank
over \(R\) is infinite.

\begin{theorem}[The polynomial diagonal action]
\label{thm:da-unit}
Evaluation at \(y=x\), followed by projection to exterior degree
zero, defines a natural even map
\[
 \pi_E:(M,Q)\longrightarrow(E(x),D(x)).
\]
It admits an \(R\)-linear even inverse up to homotopy \(i_E\)
and an odd homotopy \(h_E\) satisfying
\eqref{eq:da-chain-maps}--\eqref{eq:da-homotopy}.
They are finite compositions of polynomial divided differences.
This holds for every polynomial \(W\), every \(c\), and every
finite free matrix factorization of \(W-c\).
\end{theorem}
\begin{proof}
Eliminate the variables in the order \(y_n,y_{n-1},\ldots,y_1\).
After eliminating \(y_{i+1},\ldots,y_n\), let \(C_i\) be the
module on the remaining exterior symbols
\(\theta_1,\ldots,\theta_i\), with those eliminated variables
replaced by \(x_{i+1},\ldots,x_n\) in its differential.
Thus \(C_n=M\) and \(C_0=E(x)\).

Put \(\theta_i\) first in the exterior order to write
\(C_i=\Lambda(\theta_i)\otimes V_i\).
The tensor signs supply the corresponding reordering isomorphism.
The differential \(d_i\) on \(V_i\), which contains the earlier
exterior symbols and the boundary, has square
\[
 d_i(y_i)^2=
 \bigl(W(x_1,\ldots,x_{i-1},y_i,x_{i+1},\ldots,x_n)-c\bigr)\,1.
\]
This follows by telescoping the first \(i-1\) terms of the
potential difference. The remaining factor on \(\theta_i\)
has entries \(r_i\) and the divided difference of this displayed
potential in \(x_i,y_i\).

Apply Theorem~\ref{thm:da-one-contraction}, treating all other
coordinates as coefficients. It gives maps
\[
 \pi_i:C_i\to C_{i-1},\qquad
 i_i:C_{i-1}\to C_i,\qquad h_i:C_i\to C_i
\]
with its identities. Explicitly, its matrix \(A\) is
\((d_i(x_i)-d_i(y_i))/r_i\); its divided-difference map
evaluates only \(y_i=x_i\). Therefore
\[
 \pi=\pi_1\cdots\pi_n,\qquad i=i_n\cdots i_1,
\]
and
\begin{equation}
 h=h_n+i_nh_{n-1}\pi_n
       +i_ni_{n-1}h_{n-2}\pi_{n-1}\pi_n+\cdots
       +i_n\cdots i_2h_1\pi_2\cdots\pi_n .
 \label{eq:da-composite-h}
\end{equation}

For two consecutive contractions the formula is
\(h=h_2+i_2h_1\pi_2\). Its commutator with the differential is
\[
 (1-i_2\pi_2)+i_2(1-i_1\pi_1)\pi_2
   =1-i_2i_1\pi_1\pi_2.
\]
The chain-map equations give this equality without extra terms.
The conditions \(h_2i_2=0\), \(\pi_2h_2=0\), and
\(h_1^2=h_2^2=0\) imply \(h^2=0\); they also give
\(\pi h=hi=0\). Induction proves every claimed identity for
\eqref{eq:da-composite-h}.

The composite \(\pi\) is exactly evaluation and projection to
exterior degree zero. For a map \(f\) the tensor formula is
\(\mathcal T(f)(\omega\otimes e)=(-1)^{|f||\omega|}
\omega\otimes f(y)e\).
Projection leaves only \(|\omega|=0\), so
\eqref{eq:da-natural} holds in every number of variables.
\end{proof}

With no coordinate variables, the diagonal is the one-dimensional
even module with zero differential, and the same identities hold
with \(M=N\), \(i=\pi=1\), and \(h=0\).
For the zero boundary, all modules and maps are zero and the
identities still hold.

Every coefficient formula in this proof uses addition,
multiplication, and division by a monic coordinate difference.
It therefore holds over any commutative coefficient ring with
identity. Applying a homomorphism of coefficient rings to the
formulas preserves all the displayed identities. Flatness is
unnecessary for this particular change of coefficients: the
homotopies themselves are changed along with the objects.

\section{A mixed two-variable boundary}
\label{sec:da-mixed}

Take \(W(x,u)=x^3+xu^2\), \(c=0\), and the rank-one boundary
\[
 D_E(y,v)=
 \begin{pmatrix}0&y\\y^2+v^2&0\end{pmatrix}.
\]
Here \(r_1=x-y\), \(r_2=u-v\), and
\[
 g_1=x^2+xy+y^2+v^2,\qquad g_2=x(u+v).
\]
Order the exterior basis as \(1,\theta_1,\theta_2,\theta_2\theta_1\).
With each entry a two-by-two block, the full differential is
\[
 Q=
 \begin{pmatrix}
 D_E&r_1\,1&r_2\,1&0\\
 g_1\,1&-D_E&0&r_2\,1\\
 g_2\,1&0&-D_E&-r_1\,1\\
 0&g_2\,1&-g_1\,1&D_E
 \end{pmatrix}.
\]
Its square is \((x^3+xu^2)\,1\).

Let
\[
 A_1=\begin{pmatrix}0&1\\x+y&0\end{pmatrix},\qquad
 N=\begin{pmatrix}0&0\\u+v&0\end{pmatrix},\qquad
 P_0=\begin{pmatrix}1&0\\0&0\end{pmatrix}.
\]
The inclusion in Theorem~\ref{thm:da-unit} is
\begin{equation}
 i(e)=e+\theta_1A_1e+\theta_2Ne
                +\theta_2\theta_1(u+v)P_0e .
 \label{eq:da-mixed-inclusion}
\end{equation}
To derive the last term, first eliminate \(v\).
The matrix whose divided difference is needed is
\[
 d_2(v)=
 \begin{pmatrix}D_E(y,v)&r_1\,1\\g_1\,1&-D_E(y,v)\end{pmatrix}.
\]
Its divided difference is
\[
 \frac{d_2(u)-d_2(v)}{u-v}
   =\begin{pmatrix}N&0\\(u+v)\,1&-N\end{pmatrix}.
\]
The next inclusion, after eliminating \(y\), has components
\((e,A_1e)\). Multiplying by the displayed matrix gives
\((Ne,((u+v)\,1-NA_1)e)\), and
\((u+v)\,1-NA_1=(u+v)P_0\). This proves
\eqref{eq:da-mixed-inclusion} from the contraction formulas.

The mixed term is necessary. If it is omitted, write the resulting
map as \(i_{\mathrm{omit}}\). The coefficient of \(\theta_1\) in
\[
 Qi_{\mathrm{omit}}-i_{\mathrm{omit}}D_E(x,u)
\]
is \(-(u^2-v^2)P_0\), a nonzero polynomial matrix.
Indeed, the omitted column has only its
\(\theta_2\theta_1\) component, and the \(r_2\)-entry of \(Q\)
sends it to \(r_2(u+v)\theta_1P_0\).
Thus separate one-coordinate inclusions, with their mixed
coefficient discarded, do not give a chain map.

\section{Transporting every ordered operator product}
\label{sec:da-operators}

The contraction gives more than an equivalence of boundary objects.
Let \(B_W=\operatorname{End}_S(\Delta_W)\), with differential
\(\delta a=D_\Delta a-(-1)^{|a|}aD_\Delta\).
A homogeneous \(a\in B_W\) acts on \(M\) as
the operator \(L(a)=a\otimes1_E\).
The tensor signs imply
\[
 L(ab)=L(a)L(b),\qquad
 \partial L(a)=L(\delta a).
\]
The operator \(L(a)\) is \(R\)-linear; it acts on the infinite
polynomial pushforward module. The finite boundary algebra is
\(A_E=\operatorname{End}_R(E(x))\), with its signed commutator
differential. Projection and inclusion give a linear comparison
\[
 F_1(a)=\pi L(a)i.
\]
It commutes with differentials, but it need not preserve products
strictly. The difference is
\[
 F_1(ab)-F_1(a)F_1(b)
   =\pi L(a)(1-i\pi)L(b)i.
\]
Equation \eqref{eq:da-homotopy} expresses the middle factor by
the homotopy \(h\). Repeated products require repeated insertions
of this same \(h\).

To state every sign, the suspension \(sV\) of a parity-graded
module \(V\) has symbols \(sv\) of parity \(|v|+1\).
The tensor sign rule is
\((f\otimes g)(v\otimes w)=(-1)^{|g||v|}f(v)\otimes g(w)\).
For a differential graded algebra \(A\), form the direct sum
\[
 T^c(sA)=\bigoplus_{n\geq1}(sA)^{\otimes_R n}.
\]
The coproduct splits a word at every interior position. A
coderivation is a linear map whose coproduct is its action in
the first tensor factor plus its signed action in the second.
There is a unique coderivation \(b_A\) with components
\[
 b_{A,1}(sa)=s\,\delta a,\qquad
 b_{A,2}(sa\otimes sb)=(-1)^{|a|}s(ab),
\]
and no components of higher arity. On a word, these formulas
replace one letter by its differential or two adjacent letters
by their product, with the tensor sign rule.
The equations \(\delta^2=0\), the product rule for \(\delta\),
and associativity give \(b_A^2=0\): respectively they cancel
the one-letter, two-letter, and three-letter contributions;
operations on disjoint positions cancel with their opposite
insertion order.

An \(A_\infty\)-morphism here means a coalgebra map
\(\Phi:T^c(sA)\to T^c(sB)\) commuting with these coderivations.
Its components \(\Phi_n:(sA)^{\otimes n}\to sB\) determine
it by summing over all partitions of a word into consecutive,
nonempty blocks. A strictly unital such morphism sends \(s1\)
to \(s1\) in arity one and vanishes in higher arities whenever
an input is \(s1\).

\begin{theorem}[The boundary action at every arity]
\label{thm:da-all-arity}
For homogeneous diagonal operators \(a_1,\ldots,a_n\), define
\begin{equation}
 \Phi_n(sa_1\otimes\cdots\otimes sa_n)
   =s\bigl(\pi L(a_1)hL(a_2)h\cdots hL(a_n)i\bigr).
 \label{eq:da-operator-components}
\end{equation}
There are \(n-1\) copies of \(h\), and none when \(n=1\).
These formulas define a strictly unital \(R\)-linear
\(A_\infty\)-morphism from \(B_W\) to \(A_E\).
Its linear component sends the left scalar operator \(f(x)1\)
to \(f(x)1_E\).
\end{theorem}
\begin{proof}
Each component has parity zero on suspended words: the
unsuspended operator on the right has parity
\(\sum_j|a_j|+n-1\).
All compositions are defined on polynomial vectors, and each
word has only finitely many consecutive-block partitions.
Thus the components define a coalgebra map on the displayed
direct sum.

Expand the differential of
\(\pi L(a_1)h\cdots hL(a_n)i\) using the product rule.
Differentiating \(L(a_j)\) produces the sign
\[
 (-1)^{|a_1|+\cdots+|a_{j-1}|+j-1}.
\]
This is exactly the tensor sign for inserting \(b_{B_W,1}\)
in position \(j\) of the suspended input word.

Differentiating the copy of \(h\) between \(a_j,a_{j+1}\)
produces the sign
\[
 (-1)^{|a_1|+\cdots+|a_j|+j-1}
\]
and replaces that \(h\) by \(1-i\pi\).
The term \(1\) merges \(L(a_j)L(a_{j+1})\); its sign is the
sign for inserting \(b_{B_W,2}\) at those positions.
The term \(-i\pi\) separates the word into two factors on the
finite boundary. The first unsuspended factor has parity
\(|a_1|+\cdots+|a_j|+j-1\), so its sign is the negative of
the target component \(b_{A_E,2}(\Phi_j\otimes\Phi_{n-j})\).
Moving these terms to the other side gives exactly the
projection of \(b_{A_E}\Phi=\Phi b_{B_W}\) to words of length one.
To obtain the full identity, let
\(K=b_{A_E}\Phi-\Phi b_{B_W}\).
The coderivation rule gives
\(\Delta K=(K\otimes\Phi+\Phi\otimes K)\Delta\).
If \(K\) vanishes on shorter input words, its value on the next
word has zero coproduct. In a tensor coalgebra such an element
has length one: for each length at least two, the part of the
coproduct splitting off its first letter is injective.
The length-one projection already vanishes, so induction gives
\(K=0\) on every word.

For the unit, \(\Phi_1(s1)=s(\pi i)=s1_E\).
An initial unit in a longer word gives \(\pi h=0\);
a final unit gives \(hi=0\); an interior unit gives \(h^2=0\).
Finally, a left scalar commutes with every \(R\)-linear map
used in the contraction, which proves the scalar assertion.
\end{proof}

For \(W=z^3\) and \(E=K(z,z^2)\),
Theorems~\ref{thm:pb-rank-one} and~\ref{thm:pb-diagonal} give
\[
 H^{\bar0}(B_W)=k[x]/(x^2),\quad H^{\bar1}(B_W)=0,
 \qquad
 H^{\bar0}(A_E)=k,\quad H^{\bar1}(A_E)=k.
\]
A quasi-isomorphism would induce isomorphisms on cohomology in
each parity. These groups show that \(\Phi\) is not a
quasi-isomorphism for this boundary.

\begin{example}[A product that requires higher components]
\label{ex:da-product}
For \(W=z^3\), take
\[
 D(z)=\begin{pmatrix}0&z\\z^2&0\end{pmatrix},\qquad
 A=\begin{pmatrix}0&1\\x+y&0\end{pmatrix}.
\]
On the diagonal factorization let
\[
 \iota=\begin{pmatrix}0&1\\0&0\end{pmatrix},\qquad
 \epsilon=\begin{pmatrix}0&0\\1&0\end{pmatrix}.
\]
Both are odd, and \(\iota\epsilon\) projects to its even summand.
Directly from \(\pi\), \(i\), and \(h\),
\[
 F_1(\iota)=\begin{pmatrix}0&1\\2x&0\end{pmatrix},
 \qquad F_1(\epsilon)=0,\qquad
 F_1(\iota\epsilon)=1_E.
\]
Thus \(F_1\) is not multiplicative.
The differentials are \(\delta\iota=g1\) and
\(\delta\epsilon=r1\).
The higher component
\[
 \Phi_2(s\iota\otimes s(r1))=s1_E
\]
is nonzero: multiplication by \(r\), followed by the divided
difference in \(h\), leaves the constant vector in the
\(\theta\)-summand, and \(\iota\) sends it to the first summand.
In contrast \(\Phi_2(s\iota\otimes s\epsilon)=0\).
The arity-two equation on \((s\iota,s\epsilon)\) reads
\(0=-s1_E+s1_E\): the source product and the source differential
cancel. Dropping the nonclosed operators from the comparison
would lose this identity.
\end{example}

\section{What the boundary action detects}
\label{sec:da-detection}

In one variable, assume \(W\) is nonconstant and fix \(c\in k\).
Write \(P=W-c\), so \(P'=W'\ne0\).
For each monic divisor \(a\) of \(P\), put
\(g_a=\gcd(a,P/a)\), normalized to be monic.
Theorems~\ref{thm:pb-rank-one} and~\ref{thm:pb-diagonal}
identify the scalar cohomology map induced by
Theorem~\ref{thm:da-all-arity} with
\[
 k[x]/(W')\longrightarrow k[x]/(g_a),\qquad [f]\longmapsto[f].
\]
It is well-defined because
\(W'=a'(P/a)+a(P/a)'\) is divisible by \(g_a\).

\begin{proposition}[Joint scalar detection on a fibre]
\label{prop:da-joint-kernel}
The kernel of the map to all rank-one boundaries,
\[
 k[x]/(W')\longrightarrow
       \prod_{\substack{a\mid P\\a\text{ monic}}} k[x]/(g_a),
\]
is \(I_{\mathrm{obj}}(P)/(W')\).
These boundaries detect the scalar action on every finite free
factorization of \(P\).
\end{proposition}
\begin{proof}
The kernel consists exactly of polynomial classes whose
representatives belong to every \((g_a)\).
Their intersection is \(I_{\mathrm{obj}}(P)\) by
Theorem~\ref{thm:pb-object-ideal}.
The polynomial \(W'\) belongs to every \((g_a)\), as just
computed. The quotient therefore has the asserted meaning.
Every finite free factorization is a direct sum of the indicated
blocks by Theorem~\ref{thm:pb-all-ranks}.
\end{proof}

For \(W=x^3\) and \(c=0\), the kernel is
\((x)/(x^2)\), although evaluation gives an identity action on
every boundary. For \(W=x^2+1\) and \(c=0\), every boundary is
contractible and the entire diagonal class \(1\) has zero image.
For the same \(W\) at \(c=1\), the factorization \(K(x,x)\)
has nonzero scalar cohomology. The diagonal kernel is unchanged
by the choice of \(c\); its boundary input specifies which
fibre is being represented.

There is also a direct homotopy expressing centrality on closed
boundary maps. Let \(a\in B_W\) be closed of parity \(p\) and
let \(f\in A_E\) be closed. Put
\[
 \chi(a,f)=(-1)^{p+1}\pi L(a)h\,\mathcal T(f)i .
\]
The operator \(\mathcal T(f)=1\otimes f(y)\), with the tensor
sign rule, graded-commutes with \(L(a)\). Naturality of \(\pi\)
gives
\[
 F_1(a)f-(-1)^{p|f|}fF_1(a)
   =\pi L(a)(if-\mathcal T(f)i).
\]
Since \(\pi\mathcal T(f)i=f\), the last parenthesis is
\(-(1-i\pi)\mathcal T(f)i\). Differentiating \(\chi(a,f)\)
and using closedness gives
\[
 \partial\chi(a,f)
      =F_1(a)f-(-1)^{p|f|}fF_1(a).
\]
Thus centrality holds by an explicit homotopy.
Compatibility among these homotopies for arbitrary strings of
boundary maps is additional structure in a categorical centre;
the identity action alone does not prove a fully faithful
comparison of centres.

\section{Formal series and finite replacement}
\label{sec:da-formal}

The same construction applies to formal power series.
Replace \(R,S\) by \(k[[x]]\) and \(k[[x,y]]\), respectively,
and assume the entries of \(D(y)\) and \(W(y)\) are formal
series. The topology is specified by total degree: the
neighbourhoods of zero consist of series all of whose monomials
have total degree at least \(N\), for \(N=0,1,\ldots\).
The boundary module \(E(x)\) has the corresponding \(x\)-adic
topology.

Substitution \(y_i=x_i\) preserves total degree.
For a monomial, subtraction of its diagonal value followed by
division by \(x_i-y_i\) lowers degree by one; the polynomial
divided-difference formula proves this directly.
In each fixed total degree, a formal series has only finitely
many monomials. The same formula therefore defines a continuous
map on formal series, lowering the filtration by at most one.
Every matrix \(A_i\) in the inclusions has formal-series entries
of nonnegative degree, so each inclusion preserves the
filtration.

It follows that \(\pi\) and \(i\) preserve the filtration, and
each summand of \eqref{eq:da-composite-h} lowers it by at most
one. All identities hold coefficient by coefficient. Here the
pushforward is a complete \(R\)-module with its specified
total-degree topology, rather than the algebraic direct sum
used for polynomials. The homotopy equivalence provides the
finite complete \(R\)-module \(E(x)\) as its replacement.

At arity \(n\), if the entries of \(a_j\) have order at least
\(N_j\geq0\), the output of
\eqref{eq:da-operator-components} has order at least
\[
 \max(0,N_1+\cdots+N_n-(n-1)).
\]
Each multiplication raises order by \(N_j\), each copy of \(h\)
lowers it by at most one, and \(i,\pi\) preserve it.
This proves continuity jointly in the inputs at each fixed arity.
No infinite sum over arities is taken in
\(T^c(sB_W)\). The statement therefore gives continuous
components at every arity, without asserting convergence of
an infinite interaction series.
Changing the coefficient ring in these formal formulas is a
coefficientwise construction; no interchange of an ordinary
tensor product with an infinite product is required.

The algebraic identity action has now been constructed, including
its operator homotopies. A categorical centre comparison must
further identify the relevant coherent natural transformations
and prove the comparison faithful in the chosen category.
A cyclic pairing requires its trace and compatibility equations.
Position-dependent chiral operations require their own local
construction. These conditions specify the additional structures
to be transported beyond the boundary action.
